%TODO
\section{迭代收敛的加速方法}
\subsection{埃特金加速收敛方法}

%@see: https://mathworld.wolfram.com/AitkensDelta-SquaredProcess.html
%@see: https://en.wikipedia.org/wiki/Aitken%27s_delta-squared_process

令\begin{align*}
	x_{k+1} &\defeq \phi(x_k), \\
	x_{k+2} &\defeq \phi(x_{k+1}), \\
	y_{k+1} &\defeq x_k - \frac{(x_{k+1} - x_k)^2}{x_{k+2} - 2 x_{k+1} + x_k},
\end{align*}

\subsection{斯蒂芬森迭代法}
将不动点迭代法与埃特金加速技巧结合，
就得到了\DefineConcept{斯蒂芬森迭代法}（Steffensen's method）.

%@see: https://en.wikipedia.org/wiki/Steffensen%27s_method

\begin{align*}
	\begin{cases}
		y_k \defeq \phi(x_k), \\
		z_k \defeq \phi(y_k), \\
		x_{k+1} \defeq x_k - \frac{(y_k - x_k)^2}{z_k - 2 y_k + x_k}
	\end{cases}
	\quad(k=0,1,2,\dotsc).
\end{align*}


所有的迭代法的通用特征：
迭代时加上一个修正项，
每次迭代时该修正项不断减小，直至为零.


遇到含有指数函数的方程，可以尝试取对数.
